\section{Introduction}
\label{sec:Introduction}

\begin{table*}
\caption{Comparison among provable join-size bounds}
\centering \scalebox{0.9}{
\begin{tabular}
{@{}lccccc@{}}
\toprule[1.2pt]
Bound & \makecell[c]{Statistics}   & \makecell[c]{Cross-relation\\alignment?}    & \makecell[c]{Form}  & \makecell[c]{Storage\\per statistic} & \makecell[c]{Exact on\\stars?}          \\ \midrule
AGM~\cite{atserias2013size,atserias2008agm} & relation sizes $|R|$ & \XSolidBrush & edge-cover LP & $O(1)$ & \XSolidBrush \\
PCE~\cite{cai2019pessimistic} & sizes $+$ max degrees & \XSolidBrush & per-join rollup & $O(1)$ & \XSolidBrush \\
SafeBound~\cite{deeds2023safebound}/DSB~\cite{deeds2023degree} & (compressed) degree sequences & \XSolidBrush & functional approx. & $O(\text{segments})$ & \makecell[c]{Berge-acyclic\\only} \\
LpBound~\cite{zhang2025lpbound,abokhamis2024lp} & $\ell_p$-norms of degree seq. & \XSolidBrush & entropy LP & $O(|P|)$ & \XSolidBrush \\
CorrBound~\cite{mayer2026corrbound} & sketched inner products & \makecell[c]{\Checkmark\\(approximate)} & entropy LP & $O(\text{sketch})$ & \XSolidBrush \\
\midrule
\framework~(ours) & exact $k$-way join counts $=$ moments $\mu_\alpha$ & \makecell[c]{\Checkmark\\(exact)} & \makecell[c]{entropy LP $+$\\certificate LP} & $O(1)$ & \makecell[c]{\Checkmark\\(order $m$, or\\$m{-}1$ under\\key-collapse)} \\ \bottomrule
\end{tabular}}
\label{tab:intro_comparasion}
\end{table*}


Cardinality estimation (CE) --- predicting the output size of a query
without evaluating it --- remains, four decades on, the optimizer's
Achilles heel~\cite{leis2015howgood,han2021cardinality}.
Underestimates are the catastrophic failure mode: they steer the
optimizer toward plans designed for small data, causing intermediate
results to explode at runtime.  This observation has motivated a line
of work on \emph{pessimistic} cardinality estimation (PCE): compute a
\emph{provable upper bound} on the join size from statistics
precomputed on the base relations~\cite{cai2019pessimistic,
deeds2023degree,deeds2023safebound,abokhamis2024lp,zhang2025lpbound}.
A bound that is never violated doubles as a soundness oracle for any
estimator~\cite{stoian2026outofbounds}, and injected into an
optimizer it provably eliminates catastrophic
plans~\cite{zhang2025lpbound}.

The state of the art is LpBound~\cite{zhang2025lpbound,
abokhamis2024lp}.  Its statistics are $\ell_p$-norms of
\emph{degree sequences}: for a join attribute $X$ of relation $R$,
$\deg_R(X)$ is the multiset of value frequencies, and its norms
$\|\deg_R(X)\|_p$ feed a polymatroid linear program whose variables
are the joint entropies $h(U)$ of attribute subsets of the query
output.  Each statistic enters through a \emph{$q$-inequality}
$(1/p)h(U) + h(V\,|\,U) \le \log\|\deg_R(V\,|\,U)\|_p$ --- one
application of H\"older's inequality --- and the optimum
$2^{h(V)}$ upper-bounds $|Q|$.

\vspace{1mm}
\noindent\textbf{The blind spot: univariate statistics.}
Every statistic LpBound observes is a function of the \emph{multiset}
of degrees of one relation at a time.  But the size of a join is an
inner product: $|R\Join_X S| = \sum_x \deg_R(x)\cdot\deg_S(x)$, which
depends entirely on how the two sequences \emph{align} over the key
domain.  As a consequence, three databases can share every $\ell_p$
norm yet differ in join size by orders of magnitude:

\begin{center}\small
\begin{tabular}{l|c|c}
alignment & same degree multisets & $\sum_x a_x b_x$ \\ \hline
ranks aligned & $a{=}[8,4,2], b{=}[9,3,1]$ & $86$ \\
ranks reversed & $a{=}[8,4,2], b{=}[1,3,9]$ & $38$ \\
disjoint supports & $a{=}[8,4,2], b{=}[0,0,0]$ & $0$ \\
\end{tabular}
\end{center}

\noindent No univariate-statistic bound can separate these instances;
any tighter bound must observe \emph{joint} information across
relations.  The question is which joint statistic to buy, at what
storage cost, and with what formal guarantee.

\vspace{1mm}
\noindent\textbf{Our framework: bounds as polynomial certificates.}
We answer this question by re-founding the problem itself.  On a star
join $R_1\Join_X\cdots\Join_X R_m$, every key $x$ of the separator
$X$ carries a \emph{degree vector}
$d(x) = (d_1(x),\dots,d_m(x))\in\mathbf{N}^m$, and the join size is
the \emph{product moment} of this \emph{degree field}:
$|Q| = \sum_x \prod_i d_i(x)$.  Every statistic is an observed moment
$\mu_\alpha = \sum_x d(x)^\alpha$ --- including LpBound's, since
$\mu_{p\cdot e_i} = \|\deg_i\|_p^p$.  A bound is then nothing more
(and nothing less) than a \emph{nonnegativity certificate}: a conic
combination of observed monomials that dominates the product
polynomial pointwise, $\prod_i d_i \le \sum_\alpha c_\alpha d^\alpha$
on a worst-case domain $\Omega$.  AM--GM, Young, and H\"older's
inequalities are the closed-form points of this cone; the certificate
LP finds the best member adapted to the observed moments.  Its dual
is the \emph{alignment LP}: the largest product moment attainable by
redistributing the degree field subject to the observed moments ---
so a bound is literally a worst-case alignment, and the certificate
is its proof.

This single object reorganizes the design space:

\begin{itemize}[leftmargin=10pt,itemsep=1pt]
\item \emph{A hierarchy, not a statistic.}  Observing the moments
  $\{\alpha : |\alpha|_1\le k\}$ defines a monotone Lasserre-style
  hierarchy~\cite{lasserre2001} over the field: level 1 is the
  univariate $\ell_p$ statistics of LpBound; level 2 adds exact pair
  join counts $|R_i\Join R_j| = \mu_{e_i+e_j}$; level 3 adds triple
  counts; level $m$ is exact, since $\mu_{\mathbf{1}}$ \emph{is} the
  product moment.
\item \emph{Exactness criteria.}  We characterize when a level
  collapses to the truth: on star queries, always at order $m$, and
  at order $m{-}1$ under a \emph{key-collapse} condition (one atom is
  a key on the separator and covers the others' supports) --- both
  verified empirically on real data.
\item \emph{A computable gain predictor.}  The H\"older gap
  $\log\frac{\min\text{-conjugate}\prod\|d_i\|}{\langle d_i,d_j\rangle}$
  measures what the level-2 moment buys over level 1 per pair ---
  turning statistic selection into a principled problem.
\item \emph{A symmetric lower-bound wing.}  The same moment LP at its
  minimum is a provable \emph{lower} bound (cf.\
  xBound~\cite{stoian2026xbound}); we show where the polynomial cone
  fails there and identify the correct (hinge) basis.
\end{itemize}

\vspace{1mm}
\noindent\textbf{Instantiation and results.}
Practically, we instantiate the hierarchy inside the entropic LP:
each exact $k$-way join count yields a support constraint
$h(\cup_{i\in S}A_i) \le \log|\Join_{i\in S} R_i|$ --- the strongest
constraint obtainable from one scalar --- whose soundness is a
two-line support argument (no new inequality family is needed,
unlike~\cite{mayer2026corrbound}).  A reduced LP keeps the method
computable on queries with up to $23$ attributes, and LP duality
emits a machine-checkable proof certificate for every bound.

On the LpBound evaluation suite (JOB-light/JOB-join/JOB-range on
IMDB~\cite{leis2015howgood,leis2018query}, STATS on
StackOverflow~\cite{han2021cardinality}, DBLP dense subgraph
patterns), all with zero violations:
pair moments alone reach geomean q-error \textbf{4.97} (vs.\ 16.14)
on JOB-light, \textbf{5.05} (vs.\ 10.07) on JOB-join, and
\textbf{7.30} (vs.\ 111.0) on JOB-range; on STATS the pair arm
reaches 88.5 (vs.\ 175.3) and triples collapse it further to
\textbf{8.28} --- a $21\times$ geomean improvement, with 34\% of
queries estimated \emph{exactly}.  On DBLP the method provably ties
the baseline --- a predicted negative control where degree sequences
already determine the alignment.  The certificate LP reproduces the
entropy LP query-by-query on all 69 star queries of STATS
($\le$3\% geomean divergence), confirming that the two formulations
are views of one object.

\vspace{1mm}
\noindent\textbf{Contributions.}
\begin{itemize}[leftmargin=10pt,itemsep=1pt]
\item \textbf{Theory}: a self-contained formulation of join-size
  bounds as a moment problem on the degree field, with bounds as
  polynomial certificates and a dual alignment LP (\S\ref{sec:theory});
  a Lasserre-style hierarchy with exactness criteria (key-collapse
  lemma); and a lower-bound wing with its basis limitation.
\item \textbf{Method}: pair/triple join-count statistics as support
  constraints in the entropic LP, a reduced LP for scalability to
  23-attribute queries, and dual proof certificates
  (\S\ref{sec:method}).
\item \textbf{Evaluation}: the full LpBound suite with zero
  violations and up to $21\times$ geomean improvement, plus a
  certificate-vs-entropy cross-validation on 69 real star queries
  (\S\ref{sec:exp}).
\end{itemize}
